\documentclass[a4paper]{article}

%------------------------------------------------
% الحزم
%------------------------------------------------

\usepackage{enumitem}
\usepackage[margin=2.2cm]{geometry}
\usepackage{amsmath,amssymb}
\usepackage{graphicx}
\usepackage{array}
\usepackage{fancyhdr}
\usepackage[table]{xcolor}
\usepackage{tikz}
\usetikzlibrary{positioning,arrows.meta,shapes.geometric,calc}
\usepackage[skins]{tcolorbox}
\usepackage{polyglossia}

%------------------------------------------------
% اللغات والخطوط
%------------------------------------------------

\setdefaultlanguage[
    calendar=gregorian,
    numerals=maghrib
]{arabic}

\setotherlanguage{english}

% إذا لم تكن الخطوط مثبتة في الجهاز، تُؤخذ من المجلد fonts/
\IfFontExistsTF{Amiri}{%
    \newfontfamily\arabicfont[Script=Arabic, Scale=1.15]{Amiri}
}{%
    \newfontfamily\arabicfont[
        Script=Arabic, Scale=1.15,
        Path=fonts/, Extension=.ttf,
        UprightFont=*-Regular, BoldFont=*-Bold,
        ItalicFont=*-Italic, BoldItalicFont=*-BoldItalic
    ]{Amiri}
}

\IfFontExistsTF{AlBattar}{%
    \newfontfamily\arabicfontsf[Script=Arabic, Scale=1.2]{AlBattar}
    % خط العناوين (يعمل حتى بعد استعمال \textenglish)
    \newfontfamily\khatt[Script=Arabic, Scale=1.2]{AlBattar}
}{%
    \newfontfamily\arabicfontsf[
        Script=Arabic, Scale=1.2, Path=fonts/,
        BoldFont=ae_AlBattar.ttf
    ]{ae_AlBattar.ttf}
    % خط العناوين (يعمل حتى بعد استعمال \textenglish)
    \newfontfamily\khatt[
        Script=Arabic, Scale=1.2, Path=fonts/,
        BoldFont=ae_AlBattar.ttf
    ]{ae_AlBattar.ttf}
}

%------------------------------------------------
% الألوان: لون رئيسي هادئ + لون تنبيه فقط
%------------------------------------------------

\definecolor{asasi}{HTML}{3D6A8E}   % اللون الرئيسي: للرسوم والمخططات
\definecolor{tanbih}{HTML}{C0603A}  % لون التنبيه: للتحذيرات والحالات السلبية
\definecolor{ramadi}{HTML}{5F6B73}  % رمادي للنصوص الثانوية والأسهم
\definecolor{dakin}{HTML}{263238}   % لون العناوين

% أسماء قديمة تُحوَّل إلى اللونين فقط
\colorlet{azraq}{asasi}
\colorlet{fayruzi}{asasi}
\colorlet{banafsaji}{asasi}
\colorlet{akhdar}{asasi}
\colorlet{ahmar}{tanbih}
\colorlet{burtuqali}{asasi}

%------------------------------------------------
% الصناديق
%------------------------------------------------

\tcbset{
    asl/.style={
        enhanced,
        arc=2mm,
        boxrule=0.6pt,
        colback=white,
        colframe=black!18,
        colbacktitle=black!5,
        coltitle=dakin,
        titlerule=0pt,
        left=8pt, right=8pt, top=5pt, bottom=5pt,
        fonttitle=\khatt\bfseries,
        before skip=8pt, after skip=8pt
    }
}

% أهداف الدرس
\newtcolorbox{ahdaf}{asl, title={أهداف الدرس}}

% تعريف (خط رئيسي رفيع على الجانب)
\newtcolorbox{taarif}[1][تعريف]{asl, title={#1},
    borderline east={2pt}{0pt}{asasi}}

% مثال
\newtcolorbox{mithal}[1][مثال]{asl, title={#1}}

% تنبيه (اللون الوحيد البارز في النص)
\newtcolorbox{intabih}[1][انتبه!]{asl, colback=tanbih!5, colframe=tanbih,
    colbacktitle=tanbih, coltitle=white, title={#1}}

% ركن أستاذ اللغة العربية
\newtcolorbox{arabiyya}[1][في قسم اللغة العربية]{asl, title={#1}}

% نص أمر يكتب للأداة
\newtcolorbox{amr}[1][ramadi]{enhanced, arc=1mm, boxrule=0pt,
    colback=black!4, borderline east={2.5pt}{0pt}{#1},
    left=8pt, right=10pt, top=4pt, bottom=4pt,
    before skip=6pt, after skip=6pt}

% الخلاصة
\newtcolorbox{khulasa}{asl, title={خلاصة الدرس}}

% عنوان القسم
\newcounter{qism}
\newcommand{\qism}[1]{%
    \stepcounter{qism}%
    \par\bigskip
    {\large\khatt\bfseries{\color{asasi}\arabic{qism}.}\ {\color{dakin}#1}}\par
    \vspace{-2pt}{\color{black!20}\rule{\linewidth}{0.6pt}}\par\nopagebreak[4]
    \smallskip}

% عنوان فرعي
\newcommand{\far}[1]{\par\smallskip
    {\color{dakin}\khatt\bfseries {\color{asasi}$\blacktriangleleft$}\ #1}\par\nopagebreak[4]}

% كلمة مهمة (خط عريض فقط)
\newcommand{\kalima}[2][]{\textbf{#2}}

%------------------------------------------------
% أنماط الرسوم
%------------------------------------------------

\newcommand{\nodecolor}{}

\tikzset{
    % كتابة النص العربي من اليمين إلى اليسار داخل العقد
    % (لون النص يُعطى بالمفتاح lawn=... وليس text=...)
    every node/.append style={
        execute at begin node={\nodecolor\everypar{\beginR}\beginR},
        execute at end node=\endR},
    lawn/.code={\def\nodecolor{\color{#1}}},
    khatwa/.style={draw=#1, fill=#1!12, very thick, rounded corners=6pt,
        minimum height=1.2cm, text width=2.3cm, align=center,
        font=\khatt\bfseries},
    khatwa/.default=azraq,
    sahm/.style={-{Stealth[length=3mm]}, very thick, color=ramadi},
    raqm/.style={circle, fill=#1, lawn=white, inner sep=2pt,
        font=\small\bfseries},
}

%------------------------------------------------
% إعداد الصفحة
%------------------------------------------------

\renewcommand{\footrulewidth}{1pt}
\renewcommand{\headrulewidth}{1pt}

\rhead{\khatt المدرسة العليا للأساتذة ببوزريعة}
\lhead{\khatt ملحقة بومرداس}

\cfoot{\thepage}
\lfoot{\khatt الأستاذ: د. بن عبد الرزاق يوسف}

\parindent 0pt
\setlength{\headsep}{15pt}
\setlength{\parskip}{4pt}

\pagestyle{fancy}

\setlist[enumerate,1]{label=\arabic*), itemsep=2pt, topsep=3pt}
\setlist[itemize]{itemsep=2pt, topsep=3pt}

% الجداول
\arrayrulecolor{black!22}
\renewcommand{\arraystretch}{1.35}

% رأس جدول ملون
\newcommand{\ras}[1]{\cellcolor{black!7}\color{dakin}\bfseries #1}

% عنوان الدرس
\newcommand{\unwan}[3]{%
    \begin{center}
        {\color{asasi}\rule{\linewidth}{1.2pt}}\par\medskip
        {\khatt\color{ramadi} #1}\par\smallskip
        {\Large\khatt\bfseries\color{dakin} #2}\par\smallskip
        {\khatt\color{ramadi} #3}\par\medskip
        {\color{asasi}\rule{\linewidth}{1.2pt}}
    \end{center}}

%================================================
% بداية الوثيقة
%================================================

\begin{document}

%------------------------------------------------
% العنوان
%------------------------------------------------

\unwan{درس الفصل السادس}
{الذكاء الاصطناعي متعدد الأنماط}
{نص ــ صورة ــ صوت ــ فيديو}

\begin{ahdaf}
في نهاية هذا الدرس، تستطيع أن:
\begin{itemize}
    \item تشرح معنى «النمط» و«متعدد الأنماط».
    \item تستعمل أداة لفهم صورة أو مخطط، ثم تتحقق من إجابتها.
    \item تحوّل نصًا إلى مخطط مفاهيمي بمساعدة الأداة.
    \item تراجع المخطط وتصحح أخطاءه.
\end{itemize}
\end{ahdaf}

%================================================
\qism{ما معنى «متعدد الأنماط»؟}
%================================================

نتواصل نحن البشر بطرق كثيرة: نكتب، ونتكلم، ونرسم، ونشير.
كل طريقة من هذه الطرق نسميها \kalima{نمطًا}.

\begin{taarif}
\textbf{النمط:} شكل المعلومة: نص، أو صورة، أو صوت، أو فيديو.

\textbf{الذكاء الاصطناعي متعدد الأنماط:} أداة تفهم أكثر من نمط،
وتستطيع أن تنتج أكثر من نمط.
\end{taarif}

\begin{center}
\begin{tikzpicture}[
    namat/.style={draw=#1, fill=#1!12, thick, rounded corners=5pt,
        minimum width=2.2cm, minimum height=0.85cm, font=\khatt\bfseries}]
    % المدخلات (يمين)
    \node[font=\khatt\bfseries, lawn=ramadi] at (5,2.2) {المدخلات};
    \node[namat=akhdar]    (i1) at (5,1.2)  {نص};
    \node[namat=burtuqali] (i2) at (5,0.2)  {صورة};
    \node[namat=asasi]     (i3) at (5,-0.8) {صوت};
    \node[namat=banafsaji] (i4) at (5,-1.8) {فيديو};
    % الأداة
    \node[draw=azraq, fill=azraq!15, ultra thick, rounded corners=8pt,
        minimum width=3.2cm, minimum height=3.6cm, align=center,
        font=\khatt\bfseries] (m) at (0,-0.3) {أداة\\متعددة\\الأنماط};
    % المخرجات (يسار)
    \node[font=\khatt\bfseries, lawn=ramadi] at (-5,2.2) {المخرجات};
    \node[namat=akhdar]    (o1) at (-5,1.2)  {شرح مكتوب};
    \node[namat=fayruzi]   (o2) at (-5,0.2)  {مخطط};
    \node[namat=burtuqali] (o3) at (-5,-0.8) {صورة جديدة};
    \node[namat=asasi]     (o4) at (-5,-1.8) {قراءة صوتية};
    \foreach \i in {1,...,4} {
        \draw[sahm] (i\i.west) -- (m.east);
        \draw[sahm] (m.west) -- (o\i.east);
    }
\end{tikzpicture}
\end{center}

\begin{arabiyya}[أمثلة لأستاذ اللغة العربية]
\begin{itemize}
    \item \textbf{صورة $\leftarrow$ نص:} يصوّر صفحة مخطوطة، ويطلب كتابتها نصًا رقميًا.
    \item \textbf{نص $\leftarrow$ مخطط:} يحوّل درس «الجملة الاسمية» إلى خريطة ذهنية.
    \item \textbf{نص $\leftarrow$ صوت:} يستمع إلى قراءة لقصيدة قبل تقديمها للتلاميذ.
\end{itemize}
\end{arabiyya}

%================================================
\qism{فهم صورة باستعمال الذكاء الاصطناعي}
%================================================

يمكن أن نرفع إلى الأداة صورة فيها معلومات:
مخطط بياني، أو جدول، أو رسم توضيحي.
ثم نطلب منها شرح ما تراه.

\far{مثال على أمر}
\begin{amr}[asasi]
حلّل هذه الصورة، واشرح لي المعلومات الموجودة فيها بطريقة بسيطة.
\end{amr}

\far{خطوات العمل}
\begin{center}
\begin{tikzpicture}[node distance=7mm]
    \node[khatwa=azraq, text width=2.4cm] (a) {اختر صورة\\واضحة};
    \node[khatwa=fayruzi, text width=2.4cm, left=of a] (b) {ارفعها\\مع أمر دقيق};
    \node[khatwa=burtuqali, text width=2.4cm, left=of b] (c) {اقرأ\\الشرح};
    \node[khatwa=akhdar, text width=2.4cm, left=of c] (d) {قارن\\بالصورة};
    \draw[sahm] (a)--(b); \draw[sahm] (b)--(c); \draw[sahm] (c)--(d);
\end{tikzpicture}
\end{center}

\begin{intabih}[أين تخطئ الأداة غالبًا؟]
\begin{itemize}
    \item \textbf{الأرقام:} قد تقرأ رقمًا خطأ، أو تخلط بين عمودين.
    \item \textbf{العلاقات:} قد تفهم العلاقة بين عنصرين بشكل خاطئ.
    \item \textbf{الخط العربي:} قد تخطئ في قراءة خط اليد أو الحروف غير الواضحة.
\end{itemize}
لذلك نقارن دائمًا إجابة الأداة بالصورة الأصلية.
\end{intabih}

%================================================
\qism{تحويل النص إلى مخطط مفاهيمي}
%================================================

\begin{taarif}
\textbf{المخطط المفاهيمي:} رسم يعرض المفاهيم الأساسية لدرس ما،
ويربط بينها بأسهم توضح العلاقات.
\end{taarif}

\far{الأمر}
\begin{amr}[asasi]
حوّل النص الآتي إلى مخطط مفاهيمي بسيط،
يوضح العلاقة بين المفاهيم الرئيسية.
\end{amr}

\begin{mithal}[مثال: مخطط لدرس «الجملة الاسمية»]
\centering
\begin{tikzpicture}[
    mafhum/.style={draw=#1, fill=#1!12, thick, rounded corners=5pt,
        minimum height=0.85cm, text width=2.4cm, align=center,
        font=\khatt\bfseries},
    alaqa/.style={-{Stealth[length=2.5mm]}, thick, ramadi}]
    \node[mafhum=azraq, text width=3.2cm] (j) at (0,0) {الجملة الاسمية};
    \node[mafhum=akhdar]    (m) at (3,-1.6)  {المبتدأ};
    \node[mafhum=burtuqali] (k) at (-3,-1.6) {الخبر};
    \node[mafhum=akhdar!70!white, font=\small] (m1) at (3,-3.1) {اسم مرفوع};
    \node[mafhum=burtuqali!70!white, font=\small] (k1) at (-1.5,-3.1) {مفرد};
    \node[mafhum=burtuqali!70!white, font=\small] (k2) at (-4.5,-3.1) {جملة\\أو شبه جملة};
    \draw[alaqa] (j) -- node[above, sloped, font=\small, lawn=black] {تتكون من} (m);
    \draw[alaqa] (j) -- node[above, sloped, font=\small, lawn=black] {تتكون من} (k);
    \draw[alaqa] (m) -- node[right, font=\small, lawn=black] {هو} (m1);
    \draw[alaqa] (k) -- node[right, font=\small, lawn=black] {يكون} (k1);
    \draw[alaqa] (k) -- (k2);
\end{tikzpicture}
\end{mithal}

\far{مراجعة المخطط}
بعد أن تقترح الأداة المخطط، نسأل أنفسنا:
\begin{enumerate}
    \item هل يحتوي المخطط على \textbf{كل} الأفكار الأساسية في النص؟
    \item هل توجد علاقة \textbf{خاطئة} بين بعض المفاهيم؟
    \item هل أضافت الأداة مفهومًا \textbf{غير موجود} في النص؟
    \item ماذا أضيف أو أحذف ليصبح المخطط أوضح لتلاميذي؟
\end{enumerate}

\begin{khulasa}
\begin{itemize}
    \item \textbf{النمط} هو شكل المعلومة: نص، صورة، صوت، فيديو.
    \item الأداة متعددة الأنماط تفهم أكثر من نمط وتنتج أكثر من نمط.
    \item الأداة قد تخطئ في \textbf{الأرقام والعلاقات}، فنقارن بالأصل دائمًا.
    \item المخطط المفاهيمي يحتاج إلى \textbf{مراجعة} الأستاذ قبل استعماله في القسم.
\end{itemize}
\end{khulasa}

\end{document}
